This thesis investigates the relationship between creativity and machines through the lens of mediation theory. The research provides a thorough review of existing Theories of Concepts (TOCs) and examines their influence on the evolution of rule-based and data-driven computational approaches, in the context of creative practices. The research objective is to establish a theoretical framework able to describe the two approaches and utilize it to identify and validate critical factors of concept formation that affect the creative process, with particular focus on the observed trend towards data-driven technologies.

An extension of mediation theory is proposed, distinguishing between two computational approaches and their associated TOCs. The first two studies presented are aimed exploring the practice of concept representation using data-driven tools, with the objective of identifying the process' critical aspects. The first study is a collaboration with a music composer aimed at training a model to generate music in her unique style. The second study explores the reflective potential of dataset curation using generative adversarial networks, in partnership with a fellow researcher and photographer. The third study investigates the relationship between critical factors identified in the first two studies, within the context of text-to-image generation using StableDiffusion.

The findings highlight the significance of dataset curation for artists and designers adopting data-driven tools. From the studies, language emerges as a powerful interface for concepts, with potential implications on human and non-human creativity as large language models advance. The research indicates a possible shift in focus from product to process in creative practices, emphasizing the need for adaptation and skill development in the age of abundant content generation.